\subsection{Social Interaction and Moral Learning}
\label{sec:5.1.2}
Human moral behavior is shaped by watching other people act, being corrected, and arguing about cases. A design that treats interaction as an interface added after training has given up that channel; a design that treats it as the training has to answer for who is doing the interacting.

Two mechanisms carry most of it. Observational learning, in Bandura's sense, lets a system build a repertoire from watching someone handle real decisions — a system watching a physician allocate limited time among chronic patients is picking up a weighting that no rule list states \autocite{bandura1977social}. Evaluative feedback does what reinforcement does in human learning, strengthening some responses and weakening others, with the standing hazard that what gets strengthened is what the people giving feedback found acceptable rather than what was right. Chapter~\ref{sec:7} is what happens when that hazard is industrialized.

Emotional awareness belongs here in the form section~\ref{sec:2.2.1} settles on. What is worth engineering is compassion — steady, welfare-directed action on an accurate read of someone's state — and not manufactured shared distress, which imports a cost and buys nothing.

These channels converge on apprenticeship rather than instruction, which section~\ref{sec:5.1.3} works out.
